% GENERATED FILE. Edit canonical lessons, then run npm run generate:books.
% canonical-source-hash: fnv1a64:8d9fd394b0579a8e
% canonical-lessons: PT-C521-comedia, PT-C521-opera, PT-C521-bailado, PT-C521-bailarina, PT-C521-camarim

\chapter{Comédia, ópera e bailado}
\label{ch:pt-comedia-opera-e-bailado}
\hlchaptermodality{521}

\begin{chapteropening}
\textbf{By the end of this chapter:} \emph{I can talk about different kinds of show.}

Complete the last lesson of chapter 521: I can talk about different kinds of show.
\end{chapteropening}

\section[a comédia]{a comédia — comedy}

\label{lesson:PT-C521-comedia}



\begin{quote}
Before the new one: say the Portuguese for premiere, then the Portuguese for to premiere.
\end{quote}

\subsection*{You'll want to know: a comédia}
\textbf{a comédia} — \textquotedblleft{}comedy\textquotedblright{}.

\begin{quote}
Vimos uma comédia e rimos muito.
\end{quote}

\begin{quote}
\emph{We saw a comedy and laughed a lot.}
\end{quote}

\emph{Uma comédia romântica} — a romantic comedy.

\begin{cousinweb}
\emph{comédia} comes from Latin \emph{cōmoedia} (EtymDB), from Greek \emph{kōmōidía} \textquotedblleft{}comedy\textquotedblright{}, a word usually analysed as \emph{kômos} \textquotedblleft{}a revel\textquotedblright{} plus \emph{ōid\'{\={e}}} \textquotedblleft{}song\textquotedblright{}.
\end{cousinweb}

\begin{grammarlens}[title={What you've built}]
One more word for telling, explaining and describing at B1.
\end{grammarlens}

\subsection*{Your turn}
\begin{itemize}
\raggedright
  \item \emph{Say it:} \emph{a comédia}
  \item \emph{Say it:} \emph{a comédia}, once more
  \item \emph{Say it:} a sentence of your own with \emph{a comédia}
  \item \emph{Recall:} say the Portuguese for premiere, then the Portuguese for to premiere, then say \emph{a comédia} again
\end{itemize}

\begin{tcolorbox}[breakable,colback=teal!4,colframe=teal!35!black,title={Before you move on}]
What does a comédia mean? (\textquotedblleft{}Comedy\textquotedblright{}.) Say it once more.
\end{tcolorbox}
\section[a ópera]{a ópera — opera}

\label{lesson:PT-C521-opera}



\begin{quote}
Before the new one: say the Portuguese for to premiere, then the Portuguese for comedy.
\end{quote}

\subsection*{You'll want to know: a ópera}
\textbf{a ópera} — \textquotedblleft{}opera\textquotedblright{}.

\begin{quote}
Ela canta ópera desde criança.
\end{quote}

\begin{quote}
\emph{She has sung opera since she was a child.}
\end{quote}

The opera house is \emph{o teatro de ópera}.

\begin{cousinweb}
\emph{ópera} was borrowed from Italian \emph{opera} \textquotedblleft{}a work\textquotedblright{} (EtymDB), from Latin \emph{opera} \textquotedblleft{}work, effort\textquotedblright{}, related to \emph{opus}.
\end{cousinweb}

\begin{grammarlens}[title={What you've built}]
One more word for telling, explaining and describing at B1.
\end{grammarlens}

\subsection*{Your turn}
\begin{itemize}
\raggedright
  \item \emph{Say it:} \emph{a ópera}
  \item \emph{Say it:} \emph{a ópera}, once more
  \item \emph{Say it:} a sentence of your own with \emph{a ópera}
  \item \emph{Recall:} say the Portuguese for to premiere, then the Portuguese for comedy, then say \emph{a ópera} again
\end{itemize}

\begin{tcolorbox}[breakable,colback=teal!4,colframe=teal!35!black,title={Before you move on}]
What does a ópera mean? (\textquotedblleft{}Opera\textquotedblright{}.) Say it once more.
\end{tcolorbox}
\section[o bailado]{o bailado — ballet (a dance piece)}

\label{lesson:PT-C521-bailado}



\begin{quote}
Before the new one: say the Portuguese for comedy, then the Portuguese for opera.
\end{quote}

\subsection*{You'll want to know: o bailado}
\textbf{o bailado} — \textquotedblleft{}ballet (a dance piece)\textquotedblright{}.

\begin{quote}
Fomos ver um bailado no teatro.
\end{quote}

\begin{quote}
\emph{We went to see a ballet at the theatre.}
\end{quote}

Also \emph{o ballet}. Built on \emph{bailar} \textquotedblleft{}to dance\textquotedblright{}.

\begin{cousinweb}
\emph{bailado} is built on \emph{bailar} \textquotedblleft{}to dance\textquotedblright{}, which goes back to Late Latin \emph{ballāre} \textquotedblleft{}to dance\textquotedblright{} (Lewis \& Short).
\end{cousinweb}

\begin{grammarlens}[title={What you've built}]
One more word for telling, explaining and describing at B1.
\end{grammarlens}

\subsection*{Your turn}
\begin{itemize}
\raggedright
  \item \emph{Say it:} \emph{o bailado}
  \item \emph{Say it:} \emph{o bailado}, once more
  \item \emph{Say it:} a sentence of your own with \emph{o bailado}
  \item \emph{Recall:} say the Portuguese for comedy, then the Portuguese for opera, then say \emph{o bailado} again
\end{itemize}

\begin{tcolorbox}[breakable,colback=teal!4,colframe=teal!35!black,title={Before you move on}]
What does o bailado mean? (\textquotedblleft{}Ballet (a dance piece)\textquotedblright{}.) Say it once more.
\end{tcolorbox}
\section[a bailarina]{a bailarina — ballerina, dancer}

\label{lesson:PT-C521-bailarina}



\begin{quote}
Before the new one: say the Portuguese for opera, then the Portuguese for ballet.
\end{quote}

\subsection*{You'll want to know: a bailarina}
\textbf{a bailarina} — \textquotedblleft{}ballerina, dancer\textquotedblright{}.

\begin{quote}
A bailarina dançou sozinha no palco.
\end{quote}

\begin{quote}
\emph{The ballerina danced alone on the stage.}
\end{quote}

A man is \emph{o bailarino}.

\begin{cousinweb}
\emph{bailarina} comes from Italian \emph{ballerina} (EtymDB), built on a verb from Late Latin \emph{ballāre} \textquotedblleft{}to dance\textquotedblright{}.
\end{cousinweb}

\begin{grammarlens}[title={What you've built}]
One more word for telling, explaining and describing at B1.
\end{grammarlens}

\subsection*{Your turn}
\begin{itemize}
\raggedright
  \item \emph{Say it:} \emph{a bailarina}
  \item \emph{Say it:} \emph{a bailarina}, once more
  \item \emph{Say it:} a sentence of your own with \emph{a bailarina}
  \item \emph{Recall:} say the Portuguese for opera, then the Portuguese for ballet, then say \emph{a bailarina} again
\end{itemize}

\begin{tcolorbox}[breakable,colback=teal!4,colframe=teal!35!black,title={Before you move on}]
What does a bailarina mean? (\textquotedblleft{}Ballerina, dancer\textquotedblright{}.) Say it once more.
\end{tcolorbox}
\section[o camarim]{o camarim — (theatre) dressing room}

\label{lesson:PT-C521-camarim}



\begin{quote}
Before the new one: say the Portuguese for ballet, then the Portuguese for ballerina.
\end{quote}

\subsection*{You'll want to know: o camarim}
\textbf{o camarim} — \textquotedblleft{}(theatre) dressing room\textquotedblright{}.

\begin{quote}
Depois do espetáculo, fomos ao camarim cumprimentar a atriz.
\end{quote}

\begin{quote}
\emph{After the show, we went to the dressing room to greet the actress.}
\end{quote}

Plural \emph{camarins}.

\begin{cousinweb}
\emph{camarim} was borrowed from Italian \emph{camerino} \textquotedblleft{}a little room\textquotedblright{} (EtymDB), from Latin \emph{camera} \textquotedblleft{}a vaulted room\textquotedblright{}.
\end{cousinweb}

\begin{grammarlens}[title={What you've built}]
One more word for telling, explaining and describing at B1.
\end{grammarlens}

\subsection*{Your turn}
\begin{itemize}
\raggedright
  \item \emph{Say it:} \emph{o camarim}
  \item \emph{Say it:} \emph{o camarim}, once more
  \item \emph{Say it:} a sentence of your own with \emph{o camarim}
  \item \emph{Recall:} say the Portuguese for ballet, then the Portuguese for ballerina, then say \emph{o camarim} again
\end{itemize}

\begin{tcolorbox}[breakable,colback=teal!4,colframe=teal!35!black,title={Before you move on}]
What does o camarim mean? (\textquotedblleft{}(theatre) dressing room\textquotedblright{}.) Say it once more.
\end{tcolorbox}
